% Pregunta 293 (2019-febrero-OEP2018, n.º 3): con U = x1 + sqrt(x2), al subir los dos precios en la misma proporción la recta de balance se desplaza paralela hacia el origen; el nuevo óptimo E1 está a la misma altura x2* que E0 (la tangencia solo depende de x2 y de p2/p1) y todo el ajuste recae en x1. Autor: Víctor Gutiérrez Marcos
% Números (escala x: 2,5 cm por unidad de x1; escala y: 1,1 cm por unidad de x2):
% p1 = 2, p2 = 1, m = 4 -> x2* = (p1/2p2)^2 = 1, x1* = m/p1 - p1/(4p2) = 1,5, U0 = 2,5
% precios x1,25 (subida exagerada para que se vea) -> x2* = 1, x1* = 1,1, U1 = 2,1
% Curvas de indiferencia: x1 = U - sqrt(x2), traslaciones horizontales una de otra
\documentclass[tikz,border=6pt]{standalone}
\input{tikz-tcee}
\begin{document}
\begin{tikzpicture}[tcee]
  \ejes{$x_1$}{$x_2$}{6.2}{5.4}
  % Rectas de balance (misma pendiente -p1/p2)
  \draw[curva2] (0,4.4) -- (5,0) node[etiqueta,below,tceeazul] {$RP_0$};
  \draw[curvanueva,desplazada] (0,3.52) -- (4,0);
  \node[etiqueta,tceerojo,anchor=north east] at (0.95,2.5) {$RP_1$};
  % Curvas de indiferencia
  \draw[curva1] plot[domain=0.5:4.0,samples=80,smooth] ({2.5*(2.5-sqrt(\x))},{1.1*\x});
  \node[etiqueta,tceemorado,above] at ({2.5*(2.5-sqrt(4.0))},{1.1*4.0}) {$U_0$};
  \draw[curvanueva] plot[domain=0.5:3.0,samples=80,smooth] ({2.5*(2.1-sqrt(\x))},{1.1*\x});
  \node[etiqueta,tceerojo,left] at ({2.5*(2.1-sqrt(3.0))},{1.1*3.0}) {$U_1$};
  % Óptimos a la misma altura x2*
  \draw[guia] (3.75,1.1) -- (3.75,0) node[below] {$x_1^0$};
  \draw[guia] (2.75,1.1) -- (2.75,0) node[below] {$x_1^1$};
  \draw[guia] (0,1.1) node[left] {$x_2^*$} -- (2.75,1.1);
  \draw[flecha] (3.62,1.1) -- (2.88,1.1);
  \node[punto] at (3.75,1.1) {};
  \node[punto] at (2.75,1.1) {};
  \node[etiqueta,above right] at (3.75,1.1) {$E_0$};
  \node[etiqueta,above] at (2.78,1.12) {$E_1$};
  % Lectura
  \node[etiqueta,font=\footnotesize,align=left,tceerojo] at (4.5,4.3) {$p_1$ y $p_2$ suben igual:\\$p_2/p_1$ no cambia,\\$x_2^*$ tampoco; baja $x_1$};
\end{tikzpicture}
\end{document}
